% ============================================================
\section{Canonical Closure-Balance Realization}
\label{sec:closure-balance-realization}
% ============================================================

The counting layer of Chapter~6 determines a distinguished finite realization
of the additive horizon-system interface. The purpose of the present section is
to record that realized model explicitly: its signed-count arithmetic, its
canonical four-grade state space, its canonical boundary, and its canonical
horizon tower. These are not auxiliary choices. They are forced by the same
counting classification that selects the unique nontrivial balanced dimension.

\subsection{Canonical realized state space}

\begin{definition}[Signed-count arithmetic]\label{def:scalar-field}\lean{SignedCount}
The scalar-like quantity of the canonical realization is a \emph{signed count}:
a normalized difference of two finite tallies. It records excess pair count
over order count, or order count over pair count, without importing the usual
integer apparatus as primitive foundation.
\end{definition}

\begin{definition}[Canonical closure state space]
\label{def:closure-realized-state}
\lean{ClosureState}\lean{ClosureState.add}\lean{ClosureState.zero}\lean{closureCoordinates}\lean{stateFromClosureCoordinates}\lean{State}
The canonical closure-balance realization carries the four-grade additive state
\[
\cH_{\mathrm{cl}} := \{(x_0,x_1,x_2,x_3) : x_i \text{ is a signed count}\}.
\]
Addition, negation, and zero are defined coordinatewise. This is the
distinguished nontrivial finite state space selected by the balanced dimension
$D=4$.
\end{definition}

\begin{definition}[Canonical energy pairing]
\label{def:closure-realized-pairing}
\lean{State.pair}\lean{State.pair_self}
The canonical energy pairing on $\cH_{\mathrm{cl}}$ is the coordinate sum
\[
\langle x,y\rangle_{\mathrm{cl}}
:=
\sum_{i=0}^3 \langle x_i,y_i\rangle_{\mathrm{count}},
\]
where each coordinate pairing is induced by the signed-count energy. Its
diagonal recovers the canonical realized energy:
\[
\langle x,x\rangle_{\mathrm{cl}} = \|x\|_{\mathrm{cl}}^2.
\]
\end{definition}

\begin{theorem}[Realized state layer is explicit closure transport]
\label{thm:state-transport}
\lean{closureStateTransport}\lean{state_layer_is_closure_transport}
The theorem-bearing realized state layer introduces no additional ambient data.
It is exactly the canonical four-grade closure state transported into the state
language used throughout the manuscript. Equivalently, transporting a closure
state into the reconstructed state layer and then reading off its closure
coordinates returns the original closure state.
\end{theorem}

\begin{proof}
The transport is defined coordinatewise, and its inverse simply reads back
those same four coordinates. The round trip is therefore identity on closure
states.
\end{proof}

\subsection{Canonical full boundary and coherence tower}

\begin{definition}[True boundary]\label{def:true-boundary}\lean{trueBoundary}\lean{closureBoundary_realizes_trueBoundary}\lean{closureBoundaryState_unique}\lean{trueBoundary_unique}
The clean coherence-calculus realization carries a canonical full boundary $d$
on the four-grade state space. In closure-realized coordinates, it is exactly
the grade-lowering map that sends odd layers to the next lower even layer and
annihilates the even layers. It is defined before any horizon cut is applied.
Once that grade-lowering law is fixed, no alternative realized full boundary
remains.
\end{definition}

\begin{theorem}[Closure realization of the true boundary]\label{thm:true-boundary-realized}\lean{closureBoundary_realizes_trueBoundary}
In closure-realized coordinates, the rebuilt true boundary is exactly the
canonical grade-lowering closure boundary.
\end{theorem}

\begin{proof}
Both maps have the same explicit coordinate action: odd grades drop to the next
lower even grade, and even grades are annihilated. Hence they agree
coordinatewise.
\end{proof}

\begin{theorem}[Full boundary closes]\label{thm:full-boundary-closes}\lean{closureBoundaryState_square_zero}\lean{fullBoundary_closes}
The true boundary closes on the full coherence space:
\[
d^2=0.
\]
\end{theorem}

\begin{proof}
This is a direct calculation in the canonical four-grade boundary map: the
image of $d$ already lies in grades annihilated by one further application.
\end{proof}

\begin{definition}[Canonical coherence tower]\label{def:coherence-tower}\lean{closureProjectionState}\lean{closureProjectionFromCounting}\lean{closureProjection_realizes_coherenceProjection}\lean{closureTower}\lean{closureProjectionState_unique}\lean{closureTower_projection_unique}
The canonical coherence tower is induced by closure-grade cutoff. It reveals
the state one coordinate layer at a time. At horizon $h$, exactly the grades
strictly below $h$ are visible. Once the four-grade realization is fixed, this
tower is fixed with it, and any realized tower obeying the same grade-cut law
agrees with it pointwise.
\end{definition}

\begin{theorem}[Uniqueness of the realized tower and boundary]\label{thm:realized-uniqueness}\lean{closureProjectionState_unique}\lean{closureBoundaryState_unique}\lean{closureTower_projection_unique}\lean{trueBoundary_unique}
At the closure-state level, any admissible projection family satisfying the
forced grade-cut law agrees with the canonical closure projection, and any
admissible boundary satisfying the forced grade-lowering law agrees with the
canonical closure boundary. Consequently, at the rebuilt state level, any
realized tower or full boundary transporting those same closure laws agrees
pointwise with the canonical coherence tower and the true boundary.
\end{theorem}

\begin{proof}
The forced grade laws determine every coordinate explicitly. There is therefore
no residual choice: equality follows coordinatewise on closure states, and
transport carries that uniqueness to the rebuilt state layer.
\end{proof}

\begin{lemma}[Finite faithfulness of the canonical tower]\label{lem:coherence-tower-faithful}\lean{closureProjection_after_stable_dimension}\lean{closureShadow_after_stable_dimension}\lean{closureFaithfulTower}
For the four-grade coherence state space, the canonical tower becomes exact
after finitely many steps:
\[
\CCPi{h+4}x=x,
\qquad
\CCPibar{h+4}x=0.
\]
\end{lemma}

\begin{proof}
After four refinement steps every coordinate is visible, so no shadow remains.
\end{proof}

\begin{theorem}[Closure-Balance Realization of the Structural Interface]
\label{thm:closure-balance-realizes-interface}
\lean{closure_balance_realizes_interface}
The counting derivation of Section~\ref{sec:diophantine} determines a
distinguished realized instance of the operator interface of Part~I. Its data
are:
\begin{enumerate}[label=(\roman*), itemsep=0.2em]
\item the canonical four-grade additive state space of Definition~\ref{def:closure-realized-state},
\item the canonical energy pairing of Definition~\ref{def:closure-realized-pairing},
\item the nilpotent full boundary of Definition~\ref{def:true-boundary},
\item the canonical horizon tower of Definition~\ref{def:coherence-tower},
\item finite faithfulness of that tower.
\end{enumerate}
Consequently, the structural theorems of Part~I apply to this realized system
without further choice.
\end{theorem}

(Proof in Appendix~\ref{sec:appendix-topology}.)

\subsection{The Canonical Horizon Bundle and Horizon Complex}

\begin{definition}[Horizon bundle]
\label{def:horizon-bundle}\lean{HorizonBundle}\lean{HorizonBundle.projection}\lean{HorizonBundle.visible}\lean{HorizonBundle.shadow}\lean{HorizonBundle.pairing}\lean{HorizonBundle.pairing_self}
A \emph{horizon bundle} on the rebuilt carrier is the local state space together
with a nested family of horizon cuts. On the active four-grade realization, the
pair
\[
(\cH_{\mathrm{cl}},\{\CCPi{h}\}_{h\in\bbN})
\]
is the canonical horizon bundle: the visible part at horizon \(h\) is
\(\CCPi{h}x\), the shadow part is \(x-\CCPi{h}x\), and the bundle pairing is
the canonical state pairing \(\langle x,y\rangle_{\mathrm{cl}}\).
\end{definition}

\begin{definition}[Horizon complex]
\label{def:horizon-complex}\lean{HorizonComplex}\lean{HorizonComplex.observedBoundary}\lean{HorizonComplex.leakage}\lean{HorizonComplex.returnMap}\lean{HorizonComplex.defect}\lean{HorizonComplex.residualReturnField}
A \emph{horizon complex} is a horizon bundle together with a nilpotent full
boundary \(d\). The induced visible boundary, leakage, return, and defect on a
cut \(h\) are
\[
\CCObs{d}{h}:=\CCPi{h}d\CCPi{h},
\qquad
\CCLeak{d}{h}:=\CCPibar{h}d\CCPi{h},
\CCDefect{d}{h}:=\CCPi{h}d\CCPibar{h}d\CCPi{h}.
\]
The corresponding return map is \(\CCPi{h}d\CCPibar{h}\).
Thus the horizon complex is the native Part~I object carrying the cut/defect
calculus: full closure holds before the cut, and failure of closure appears
only after observation.
\end{definition}

\begin{theorem}[Closure Balance already yields the canonical horizon complex]
\label{thm:closure-balance-yields-the-canonical-horizon-complex}\lean{canonicalHorizonBundle}\lean{canonicalHorizonComplex}\lean{canonicalHorizonComplex_surface}\lean{canonical_visible_after_stable_dimension}\lean{canonical_shadow_after_stable_dimension}\lean{HorizonComplex.returnedResidual_eq_schur}
Closure Balance determines, on the rebuilt four-grade carrier, a canonical
horizon bundle and a canonical horizon complex. Concretely:
\begin{enumerate}[label=(\roman*), leftmargin=2em, itemsep=0.2em]
\item the full boundary on that complex squares to zero;
\item on every horizon cut, the observed boundary satisfies
      \[
      \CCObs{d}{h}^2=-\,\CCDefect{d}{h};
      \]
\item after four further refinements the visible cut becomes exact and the
      shadow vanishes;
\item once a shadow propagator is chosen on a cut, the corresponding returned
      residual field is exactly the Schur correction across that cut.
\end{enumerate}
So the rebuilt state carrier is not merely a four-coordinate list. It already
carries one canonical horizon complex with its native boundary and
residual-return grammar.
\end{theorem}

\begin{proof}
Take the canonical horizon bundle to be the rebuilt state carrier with the
canonical coherence tower, and take the canonical horizon complex to be that
same bundle with the true boundary \(d\). Full closure is
Theorem~\ref{thm:full-boundary-closes}, the cut-defect law is
Corollary~\ref{cor:HFT1-boundary}, and finite exact visibility is
Lemma~\ref{lem:coherence-tower-faithful}. The final clause is the Schur
factorization of the residual-return field on that same cut.
\end{proof}

\subsection{Canonical HFT Consequences}

\begin{corollary}[Closure-Balance HFT-1]
\label{cor:HFT1-boundary}\lean{closureBoundaryState_square_zero}\lean{fullBoundary_closes}\lean{closureBoundary_closes}\lean{HFT1_closure_exact}\lean{closureObservedBoundaryState_square_eq_neg_defect}
Let $\cH$ be the full coherence space, let $d$ denote its true boundary, and
let $(\CCPi{h})_{h\in\bbN}$ be the canonical coherence tower. Then $d^2 = 0$ in
the full space, observation is the cut $\CCObs{d}{h} := \CCPi{h} d \CCPi{h}$,
and
\[
\CCObs{d}{h}^2 \;=\; -\,\CCDefectExpr{d}{h}.
\]
Thus closure holds in the full space and fails only after the horizon cut.
\end{corollary}

\begin{proof}
Apply Theorem~\ref{thm:HFT1-general} to the true boundary $d$ furnished by
Theorem~\ref{thm:full-boundary-closes}.
\end{proof}

\begin{corollary}[Closure-Balance HFT-2]
\label{thm:HFT2-canonical}\lean{closureEnergyTower}\lean{closureFaithfulTower}\lean{HFT2_closure_finite}\lean{HFT2_closure_stabilized}\lean{closureHFT2_state_finite}\lean{closureHFT2_state_stabilized}
For the canonical coherence tower, adjacent horizons close internally and the
refinement layers are exact coherence increments. Hence for nested coherence
horizons $h_0 \le h_1 \le \cdots \le h_N$ and any $x\in\cH$,
\[
\|x - \CCPi{h_0} x\|^2
=
\sum_{j=0}^{N-1} \|(\CCPi{h_{j+1}} - \CCPi{h_j})x\|^2 + \|x - \CCPi{h_N} x\|^2.
\]
Since the canonical coherence tower is faithful, the tail vanishes after
sufficient refinement. Hence there exists $N$ such that for every $k \ge 0$,
\[
\|x - \CCPi{h_0} x\|^2
=
\sum_{j=0}^{N+k-1} \|(\CCPi{h_0+j+1} - \CCPi{h_0+j})x\|^2.
\]
\end{corollary}

\begin{proof}
Apply Theorem~\ref{thm:HFT2-general} to the canonical coherence tower supplied
by Definition~\ref{def:coherence-tower} and
Lemma~\ref{lem:coherence-tower-faithful}. In this realized system, the tower
hypotheses are not postulated externally: they are forced by Closure Balance.
\end{proof}


% --- Hodge scalar and canonical invariants ---

\subsection{The Hodge scalar identity on the balanced complex}
\label{subsec:hodge-scalar-theorem}

The six-slot carrier of the balanced four-incidence is simultaneously the edge
carrier of the clique complex \(K_4\). At that level, the simplicial Hodge
\(1\)-Laplacian is maximally symmetric.

\paragraph{Hodge scalar identity on \(K_4\).}
Let
\[
\Delta_1:=\partial_1^\top\partial_1+\partial_2\partial_2^\top
\]
denote the simplicial Hodge \(1\)-Laplacian on the edge space of \(K_4\).
Then
\[
\Delta_1 = 4 I_6.
\]
The proof is the same exact cancellation mechanism used in the short-paper
companion treatment:
every diagonal entry contributes \(2\) from shared vertices and \(2\) from
shared triangles, while every adjacent off-diagonal contribution is
\(+1\) from the lower term and \(-1\) from the upper term, hence cancels.
Disjoint edges contribute \(0\) to both pieces. Thus no edge direction is
preferred on the balanced complex.

\paragraph{Why this matters here.}
At the discrete balanced level, nilpotent closure therefore distributes its
content democratically across the six-slot carrier. The later channel
distinction is not already visible in the raw Hodge operator on \(K_4\); it
appears only after the intrinsic \(1+2+3\) decomposition is read through the
continuum bridge. In that sense the balanced complex is maximally symmetric
before the continuum carrier readouts differentiate the phase, wave, gauge,
and geometric lanes.

\subsection{Canonical horizon invariants}

Two a priori different horizon-dependent data packages are now available.
Chapter~5 attached a shadow-transport spectrum, its moments, and its
determinant polynomial to the boundary leakage operator. The canonical tower
also attaches a tail-energy profile to each state. In the closure-balance
realization these behave very differently: the boundary package collapses
identically, whereas the state-tail package yields a nontrivial finite profile.

\begin{theorem}[Canonical collapse of boundary shadow transport]
\label{thm:canonical-shadow-collapse}
\lean{canonical_shadow_collapse}
For the canonical true boundary $d$ and the canonical coherence tower
$(\CCPi{h})_{h\in\bbN}$,
\[
\CCPibar{h} d \CCPi{h} = 0,
\qquad
\CCDefectExpr{d}{h}=0
\]
for every horizon $h$. Consequently, on the finite observed sector
$\CCPi{h}\cH_{\mathrm{cl}}$, the shadow-transport invariant family of
Definition~\ref{def:shadow-transport-invariants} is trivial:
\[
\Lambda_h(d)=\{0,\dots,0\},
\qquad
M_k(h)=0 \quad (k\ge 1),
\qquad
\Delta_h(t)=1.
\]
\end{theorem}

\begin{proof}
In closure-realized coordinates, $\CCPi{h}$ keeps exactly the grades strictly
below $h$, while $d$ lowers each odd grade to the next lower even grade and
annihilates the even grades. Hence if $x$ is observable at horizon $h$, every
nonzero coordinate of $dx$ still lies in a grade strictly below $h$. Therefore
$\CCPibar{h} d \CCPi{h}x=0$ for every $x$, proving $\CCPibar{h} d \CCPi{h}=0$.
The defect operator $\CCPi{h} d \CCPibar{h} d \CCPi{h}$ therefore vanishes as
well. The statements about $\Lambda_h(d)$, the moments, and the determinant
polynomial then follow from Theorem~\ref{thm:shadow-transport-package}.
\end{proof}

\begin{definition}[Canonical tail and revelation profiles]
\label{def:canonical-tail-profile}
\lean{canonicalGradeEnergy}\lean{canonicalTailProfile}\lean{canonicalRevelationProfile}
For a closure-realized state $x=(x_0,x_1,x_2,x_3)\in\cH_{\mathrm{cl}}$, define
its \emph{grade energies}
\[
\varepsilon_i(x) := \|x_i\|_{\mathrm{count}}^2
\qquad (0\le i\le 3).
\]
Define its \emph{tail profile} by
\[
E_h(x) := \|x-\CCPi{h}x\|_{\mathrm{cl}}^2,
\]
and its \emph{revelation profile} by the one-step increment energies
\[
\mu_h(x) := \|(\CCPi{h+1}-\CCPi{h})x\|_{\mathrm{cl}}^2
\qquad (h\in\bbN).
\]
\end{definition}

\begin{theorem}[Exact canonical tail and revelation laws]
\label{thm:canonical-tail-laws}
\lean{canonical_tail_laws}
For every closure-realized state $x=(x_0,x_1,x_2,x_3)$,
\begin{align*}
E_0(x) &= \varepsilon_0(x)+\varepsilon_1(x)+\varepsilon_2(x)+\varepsilon_3(x),\\
E_1(x) &= \varepsilon_1(x)+\varepsilon_2(x)+\varepsilon_3(x),\\
E_2(x) &= \varepsilon_2(x)+\varepsilon_3(x),\\
E_3(x) &= \varepsilon_3(x),\\
E_h(x) &= 0 \qquad (h\ge 4),
\end{align*}
and
\begin{align*}
\mu_0(x) &= \varepsilon_0(x),\\
\mu_1(x) &= \varepsilon_1(x),\\
\mu_2(x) &= \varepsilon_2(x),\\
\mu_3(x) &= \varepsilon_3(x),\\
\mu_h(x) &= 0 \qquad (h\ge 4).
\end{align*}
\end{theorem}

(Proof in Appendix~\ref{sec:appendix-topology}.)

\begin{corollary}[Finite cumulative law and exact reconstruction]
\label{cor:canonical-tail-reconstruction}
\lean{canonical_tail_reconstruction}
For every closure-realized state $x$ and every horizon $0\le h\le 3$,
\[
E_h(x)=\sum_{j=h}^{3}\mu_j(x).
\]
Moreover $E_h(x)=0$ for every $h\ge 4$.
Equivalently,
\[
E_0(x)=\mu_0(x)+E_1(x),\quad
E_1(x)=\mu_1(x)+E_2(x),\quad
E_2(x)=\mu_2(x)+E_3(x),\quad
E_3(x)=\mu_3(x).
\]
Consequently,
\[
E_0(x)\ge E_1(x)\ge E_2(x)\ge E_3(x)\ge E_4(x)=0,
\]
and the grade energies are recovered exactly from the horizon profile:
\[
\varepsilon_i(x)=\mu_i(x)
\qquad (0\le i\le 3).
\]
\end{corollary}

\begin{proof}
Substitute the formulas of Theorem~\ref{thm:canonical-tail-laws}. The
cumulative identity is immediate, the monotonicity follows because each
$\mu_i(x)$ is a nonnegative energy, and the final identity
$\varepsilon_i(x)=\mu_i(x)$ is exactly the one-step revelation law.
\end{proof}

\paragraph{What varies canonically.}
The closure-balance realization therefore separates two notions that might
otherwise be conflated. The true boundary is internally exact across every
horizon cut, so its shadow-transport spectrum is trivial. The nontrivial
horizon dependence lies instead in the state-side tail and revelation
profiles: the horizon reveals one grade at a time, and the successive energy
drops record exactly how much of the state occupied each hidden grade. This is
a purely internal exposure law of the canonical tower; no external scale
parameter is assumed here.

\paragraph{Distinguished realized model.}
Chapter~6 does not enlarge the structural calculus. It identifies the canonical
finite model singled out by the counting bedrock and shows that the structural
operator identities of Part~I apply to that model without any further choice.
